\chapter{Как машина учится двигаться}
% полная версия: chapters/15_motorlearning.tex

Мы видели два способа учиться. Без учителя --- «вместе, значит связаны» --- машина запоминает то, что часто встречается. С дофамином она учится на сюрпризах: вышло лучше или хуже, чем ожидалось. Но когда вы промахнулись, бросая мяч в корзину, вы знаете не только, что вышло плохо. Вы знаете, \emph{насколько} и \emph{в какую сторону} промахнулись. Это третий учитель --- ошибка с направлением, личная для каждого выхода.

\section*{Провод ошибки}

В мозге есть отдельная схема для обучения движениям, и устроена она очень необычно. Огромное число мелких \nt{ГЛУТ}-нейронов, десятки миллиардов, раскладывает входные сигналы в очень длинный список признаков. Выходные нейроны этой схемы --- тормозные, \nt{ГАМК}, их десятки миллионов, и каждый слушает сотню тысяч входов. И к каждому выходному нейрону тянется ровно один особый провод --- провод ошибки. Если ошибка пришла, пока вход был активен, этот вход ослабевает. Это классическое правило «уменьшай ошибку», которым инженеры учат адаптивные фильтры.

Такой учитель быстрее дофаминового: у каждого выхода своя ошибка, и не нужно выуживать свою долю из общего сигнала на всех. Правда, ошибка приходит с опозданием, и соединению снова нужен «стикер», помнящий, что оно участвовало. Опыты показывают, что длительность этого стикера подобрана под задержку ошибки.

\section*{Внутренняя модель}

Что выучивает такая схема? Модель собственного тела: какой результат даст команда. С ней не нужно ждать запаздывающих датчиков --- результат можно предсказать заранее. По одной из гипотез, поэтому вы не можете пощекотать сами себя: мозг заранее знает, что будет, и вычитает предсказанное.

\section*{Быстрый и медленный ученик}

Наденьте очки, сдвигающие картинку вбок, и бросайте дротики. Первые броски уходят в сторону, но за десятки попыток вы приспособитесь. Снимите очки --- и снова промахи, теперь в другую сторону. Интересное начинается дальше. Опыты хорошо описываются двумя учениками сразу: быстрый быстро учится и быстро забывает, медленный учится медленно, но помнит долго. Если после долгого приспособления коротко переучиться обратно, общая поправка вернётся к нулю --- но старое умение затем само частично всплывёт: быстрый ученик забыл, а медленный помнит. Модель с одним учеником так не умеет, а люди --- умеют.

Зачем два ученика? По гипотезе, это тот же рецепт, что у штурмана из девятой главы: мир меняется и быстро, и медленно, и разумно следить за обоими.

\fullversion{15}
